% ---------------------------------------------------------------------
%  Chapter 3 : Specifying a Model --- the TheoryBuilder
% ---------------------------------------------------------------------
\chapter{Specifying a Model: the \texttt{TheoryBuilder}}\label{ch:theory-spec}

In the previous chapter you loaded a theory with a single call and ran it. That
theory came from a small file under \file{theories/}, and the heart of that file
was a chain of method calls ending in \code{.build()}. This chapter opens that
chain up. By the end you will know what every method does, why the model you
describe in a few readable lines becomes a large Python dictionary full of
symbolic functions, and exactly where the bodies of those functions come from.

This is the \emph{front door} of the whole pipeline. Everything downstream ---
the propagator, the mean-field solve, the diagram enumeration, the loop integrals
--- consumes one object, the \term{model dict}, and this subsystem is the only
thing that produces it. It never integrates a diagram or solves an equation; its
single job is \emph{input specification and compilation}. Understand it well and
the rest of the manual is the story of what happens to its output.

\begin{note}
The three source files this chapter documents are
\file{pipeline/theory.py} (the builder classes and \code{build()}),
\file{pipeline/theory\_compiler.py} (which turns the text you type into Sage
lambda functions), and \file{pipeline/theory\_templates.py} (pre-canned actions
and kernels). Wherever a claim here can be checked against the code, the
\code{file:line} is given so you can read the original.
\end{note}

\section{Why a builder at all?}

To compute cumulants, the rest of the pipeline needs a \term{model dict}: a large
Python \code{dict} whose values are not merely numbers and metadata but several
\emph{Sage-aware lambda functions}. The keys with names like \code{action},
\code{mf\_bg\_conditions}, \code{kernel\_ft\_image}, \code{phi\_concrete}, and
\code{mf\_equations} all map to callables. Each such lambda takes a single runtime
argument conventionally called \code{ns} --- a ``namespace'' object that carries
the symbolic field and parameter variables --- and returns a \emph{SageMath
symbolic expression}.

\begin{defn}
\textbf{SageMath} is a large open-source mathematics system built on Python. Its
relevant feature here is a \emph{computer algebra system}: it can hold and
manipulate \emph{symbolic} expressions (variables, functions, integrals) rather
than only floating-point numbers. The central object is the \term{Symbolic Ring},
abbreviated \code{SR}. \code{SR.var('mu')} creates a symbolic variable \code{mu};
arithmetic on \code{SR} objects builds expression trees, so
\code{SR.var('a')*SR.var('x')\^{}2} is the \emph{symbolic} $a\,x^{2}$, not a
number. Sage can differentiate, Taylor-expand, Fourier-transform, and substitute
on these trees.
\end{defn}

Writing one of these dicts by hand means writing those lambdas by hand, which
demands fluency in both SageMath's symbolic algebra and a thicket of private
naming conventions (which internal name is the response field, which is the
saddle, how a transfer function's derivatives are named, and so on). The module
docstring of \file{pipeline/theory.py:1} opens by saying exactly this: the
pipeline ``takes a Python dict (\code{HAWKES\_MODEL}) with several Sage-aware
lambdas \ldots Writing one of these by hand requires familiarity with Sage's
symbolic algebra and the framework's conventions.''

The builder is the human-friendly alternative. Instead of hand-writing the dict,
you \emph{describe} the theory with a fluent chain of method calls --- declare a
field, declare a parameter, type the action as a string, type the mean-field
equation as a string --- and call \code{.build()}. The builder accumulates your
declarations, compiles the text strings into the Sage lambdas the pipeline
expects, and assembles the dict for you.

\begin{gotcha}
The top docstring of \file{pipeline/theory.py} (lines~1--62) still reads like a
roadmap for a half-finished prototype: it lists an older API
(\code{transfer\_function}, \code{use\_action\_template}, \code{set\_action}) and
describes text compilation as future work. Much of that roadmap is now done.
\emph{Do not trust the top docstring over the method bodies} --- this chapter
follows the bodies.
\end{gotcha}

\section{Choosing a builder class}

There are three concrete builder classes, all defined near the end of
\file{pipeline/theory.py}:

\begin{description}
  \item[\code{TemporalTheoryBuilder}] (\file{pipeline/theory.py:2087}) ---
        time-only theories: a scalar variable $x(t)$, or several coupled
        time-only variables, with no spatial extent. It carries the
        temporal-only methods (populations, convolution kernels, coloured-noise
        embedding). Its \code{build()} asserts that \emph{every} field is
        non-spatial.
  \item[\code{SpatialTheoryBuilder}] (\file{pipeline/theory.py:2102}) --- spatial
        PDE / field theories, a field $\varphi(x,t)$. It carries the spatial-only
        methods (\code{spatial\_dim}, \code{boundary}, \code{initial},
        \code{dyson}, \code{reference\_diffusion}). Its \code{build()} asserts
        that \emph{at least one} field is spatial. Note its constructor defaults
        \code{n\_populations=0} (\file{pipeline/theory.py:2108}).
  \item[\code{TheoryBuilder}] (\file{pipeline/theory.py:2121}) --- a
        back-compatibility class that mixes in \emph{both} method sets and
        auto-detects spatial-versus-temporal at \code{build()} time. New code
        should prefer the two specific classes.
\end{description}

The split is done with Python \emph{mixins} --- small classes that contribute
methods. \code{TemporalTheoryBuilder} inherits from \code{\_TemporalMethods} and a
shared base \code{\_BaseTheoryBuilder}; \code{SpatialTheoryBuilder} from
\code{\_SpatialMethods} and the same base; \code{TheoryBuilder} from all three.
The auto-detect lives in a method each subclass overrides,
\code{\_resolve\_spatial\_dim}. The temporal subclass raises if any field declared
a non-zero spatial dimension; the spatial subclass raises if none did
(\file{pipeline/theory.py:2092}, \code{:2111}):

\begin{lstlisting}
class TemporalTheoryBuilder(_TemporalMethods, _BaseTheoryBuilder):
    def _resolve_spatial_dim(self) -> int:
        d = super()._resolve_spatial_dim()
        if d > 0:
            raise ValueError(...'Use SpatialTheoryBuilder...')
        return 0
\end{lstlisting}

Every builder method returns \code{self}, which is what makes the fluent
chain-of-dots style work: each call hands the builder back so the next dot can
land on it.

\section{Declaring fields}

\msrjd{} field theory doubles every observable field $\varphi$ with a
\term{response field} $\tilde\varphi$ (the ``tilde'' field). The physics is
recalled in the next section; for now the point is that declaring one physical
field implies a small family of related symbols, and the builder generates that
family for you.

\subsection{\texttt{physical\_field}: one call, three symbols}

\code{physical\_field} (\file{pipeline/theory.py:806}) is the single most
important field method. It supports two calling styles, distinguished by whether
you pass \code{natural\_name}:

\begin{description}
  \item[New style] (\code{natural\_name} omitted) --- the \code{name} you pass
        \emph{is} the natural letter. \code{.physical\_field('n')} means: the
        user-facing letter is \code{n}, and the internal \term{fluctuation field}
        is derived as \code{dn}. So you write \code{n} (or \code{dn}; both refer
        to the fluctuation) and the framework knows the internal name is
        \code{dn} (\file{pipeline/theory.py:858}).
  \item[Legacy style] (\code{natural\_name} given) --- \code{name} is the literal
        internal name and \code{natural\_name} is the separate user letter, as in
        \code{.physical\_field('dn', natural\_name='n')}. This is how the older
        hand-written theory files read.
\end{description}

Whichever style you use, two further symbols are generated automatically unless
you opt out:

\begin{itemize}
  \item \textbf{Auto-response} (unless \code{auto\_response=False}): a conjugate
        response field named \code{<natural>t}. So \code{.physical\_field('n')}
        also declares \code{nt} (\file{pipeline/theory.py:885}). It is skipped if
        a response field of that name already exists.
  \item \textbf{Auto-saddle} (unless \code{auto\_saddle=False}): a saddle-point
        parameter named \code{<natural>star} carrying \code{mean\_field=True}, so
        \code{.physical\_field('n')} also declares \code{nstar}
        (\file{pipeline/theory.py:899}). Its default domain is \code{'positive'}
        for the rate field \code{n} and free for everything else.
\end{itemize}

So the single line \code{.physical\_field('x')} quietly creates three things: the
fluctuation field \code{dx}, the response field \code{xt}, and the saddle
parameter \code{xstar}. The next two sections explain why those three exist.

\begin{defn}
A \textbf{saddle} (or \textbf{mean field}) is the deterministic steady state
$\bar\varphi$ around which the theory is expanded; internally it is the
\code{*star} parameter (\code{xstar}, \code{nstar}, \ldots). It is not a number
you supply --- the pipeline solves for it. The \textbf{fluctuation field} is
$\delta\varphi = \varphi - \bar\varphi$, the internal field the action is
Taylor-expanded in (\code{dx}, \code{dn}, \ldots).
\end{defn}

For a spatial field, pass \code{spatial\_dim=d}. The dimension is a
\emph{per-field} property by design (\file{pipeline/theory.py:79}); the
bulk-setter \code{.spatial\_dim(d)} (\file{pipeline/theory.py:122}) just stamps it
onto every field at once and becomes the default for fields declared afterwards.
Values $d \in \{1,2,3\}$ are validated end-to-end; $d=0$ is time-only. All spatial
fields in one theory must share a single non-zero dimension (a v1 constraint
enforced in \code{build()}).

\code{response\_field} (\file{pipeline/theory.py:796}) declares a response field
directly, which you rarely need given auto-response, but it is the escape hatch
when a response field has no physical partner.

\section{Declaring parameters}

\code{parameter} (\file{pipeline/theory.py:914}) records one model parameter as a
\code{ParameterSpec}. The signature and its record type are:

\begin{lstlisting}
@dataclass
class ParameterSpec:
    name: str
    indexed: bool = False         # True = vector
    matrix: bool = False          # True = N x N matrix
    domain: Optional[str] = None  # e.g. 'positive'
    default: Any = None
    mean_field: bool = False      # True for saddle quantities (nstar, ...)
    natural_name: Optional[str] = None
    indexed_by: Optional[list[str]] = None  # [] scalar, ['E'] vec, ['E','I'] mat
\end{lstlisting}

The fields that matter day-to-day are:

\begin{description}
  \item[\code{default}] the numeric value used when a run does not override it
        (scalar, list, or list-of-lists).
  \item[\code{domain}] an assumption Sage's simplifier and Fourier transform can
        exploit, most commonly \code{'positive'}. Passing
        \code{domain='positive'} creates the variable as
        \code{SR.var(name, domain='positive')}, which lets Sage's underlying CAS
        dispatch the right simplification.
  \item[\code{mean\_field}] \code{True} flags the parameter as a saddle quantity
        (the \code{*star} variables). Auto-saddle sets this for you.
  \item[\code{indexed\_by}] the modern heterogeneous-population annotation:
        \code{[]} is a scalar, \code{['E']} a vector, \code{['E','I']} a matrix.
        The legacy flags \code{indexed} / \code{matrix} say the same thing; both
        are normalised at declaration time into an internal
        \code{(is\_vec, is\_mat)} pair.
\end{description}

A \emph{matrix} parameter (say a connectivity \code{w}) becomes, at build time, a
grid of \code{SR} variables \code{w11, w12, \ldots}, each carrying the declared
domain so the CAS can reason about them (\file{pipeline/theory.py:1634}).

\section{Writing the action as text}

The action is the centrepiece of the specification, and it is typed as a string.
Before the syntax, a one-page recap of what the action \emph{is}.

\subsection{The \msrjd{} action in one page}

A stochastic equation of motion for a field $\varphi$ (a scalar $x(t)$ or a field
$\varphi(x,t)$) driven by noise $\eta$,
\begin{equation}
  \Dt\varphi \;=\; F[\varphi] + \eta,
  \qquad \avg{\eta(t)\,\eta(t')} = 2D\,\delta(t-t'),
  \label{eq:eom}
\end{equation}
is recast as a path integral by introducing the response field $\tilde\varphi$.
The partition function becomes $Z = \int\!\mathcal{D}\varphi\,\mathcal{D}\tilde\varphi\,
\ee^{-S[\varphi,\tilde\varphi]}$ with the \msrjd{} action
\begin{equation}
  S[\varphi,\tilde\varphi] \;=\; \int\!\dd t\;
  \Big[\, \tilde\varphi\,\big(\Dt\varphi - F[\varphi]\big)
        \;-\; D\,\tilde\varphi^{2} \,\Big].
  \label{eq:action}
\end{equation}
The term $\tilde\varphi\,\Dt\varphi$ carries the causal/kinetic structure,
$\tilde\varphi\,F[\varphi]$ the drift, and $-D\,\tilde\varphi^{2}$ encodes the
white Gaussian noise as a \emph{second cumulant} on the response field. This last
fact generalises: non-Gaussian white noise of $n$-th cumulant $\kappa^{(n)}$ adds
a response-field monomial of degree $n$, namely
$-(\kappa^{(n)}/n!)\,\tilde\varphi^{n}$. That is exactly why
\code{set\_action\_text}'s docstring states that ``non-Gaussian white noise is a
response-field monomial of degree $\ge 3$'' (\file{pipeline/theory.py:1040}).

\subsection{The saddle--fluctuation split}

The cumulants are computed perturbatively around the deterministic steady state,
the saddle $\bar\varphi$. Every physical field is split
$\varphi = \bar\varphi + \delta\varphi$, and the action is Taylor-expanded in the
fluctuation $\delta\varphi$. The expansion produces a constant term, a linear
(\term{tadpole}) term, a quadratic (propagator) term, and cubic-and-higher
(interaction-\term{vertex}) terms.

\begin{defn}
The \textbf{saddle equation} is the condition that the linear tadpole term
vanishes --- physically, that $\bar\varphi$ is the deterministic fixed point. For
the quartic OU model \eqref{eq:ouquartic-here} below it reads
$(\Dt+\mu)\bar x = -\varepsilon\bar x^{3}$, which at steady state ($\Dt\to 0$) is
$\mu\bar x = -\varepsilon\bar x^{3}$.
\end{defn}

\subsection{The text DSL}

\code{set\_action\_text} (\file{pipeline/theory.py:1008}) stores the action as a
Sage-syntax string; its long docstring (lines~1008--1049) is the canonical
reference. The rules are:

\begin{itemize}
  \item You type the \emph{per-population integrand} $S_i$. The index \code{i} is
        an implicit free index; at build time the integrand is summed over
        \code{i in range(n\_populations)}. For a scalar theory that sum has one
        term.
  \item Inner sums (e.g.\ over a connectivity) use Python comprehension syntax:
        \code{sum(w[i, j] * g * dn[j] for j in pop)}, where
        \code{pop = range(n\_populations)} is pre-bound.
  \item \code{Dt} is an inert symbol standing for $\Dt$, used multiplicatively
        (\code{(Dt+mu)*x[i]}).
  \item The response field is the \code{<letter>t} name (\code{xt}, \code{nt}).
  \item You may write the \emph{full physical field} \code{x[i]} and the framework
        expands it to \code{xstar[i] + dx[i]} for you, or write the fluctuation
        \code{dx[i]} explicitly. (The mechanism behind the convenient \code{x[i]}
        form is the \code{\_FullPhysicalField} wrapper, covered in
        \S\ref{sec:textcompile}.)
\end{itemize}

The quartic OU action of the previous chapter,
\begin{equation}
  S \;=\; \int\!\dd t\;\Big[\, \tilde x\,\big((\Dt+\mu)\,x + \varepsilon x^{3}\big)
        \;-\; D\,\tilde x^{2} \,\Big],
  \label{eq:ouquartic-here}
\end{equation}
is typed exactly as it reads (\file{theories/ou\_quartic.theory.py}):

\begin{lstlisting}
.set_action_text('sum(xt[i]*((Dt+mu)*x[i] + eps*x[i]^3) - D*xt[i]^2 for i in pop)')
\end{lstlisting}

Here \code{xt} is $\tilde x$, the \code{(Dt+mu)*x[i] + eps*x[i]\^{}3} piece is the
deterministic dynamics (with the cubic \code{eps*x[i]\^{}3} the single interaction
vertex), and \code{- D*xt[i]\^{}2} is the white-noise vertex. Coupled multi-field
theories carry one such term per field, with cross-coupling appearing as
off-diagonal terms.

\section{Functions and the transfer function}

Some drifts contain a nonlinear function whose \emph{Taylor coefficients at the
saddle} are the interaction couplings. The archetype is a neural \term{transfer
function} $\phi(v)$ (the firing-rate nonlinearity). Note the unfortunate name
clash: in this codebase the transfer function is the user-declared function
literally named \code{phi}, distinct from the field $\varphi$.

\code{define\_function} (\file{pipeline/theory.py:969}) records a function by name,
formal argument list, and a Sage-syntax body, e.g.\
\code{.define\_function('phi', ['v'], 'a*v\^{}2')}. The function body may reference
any declared parameter. \code{set\_transfer\_function}
(\file{pipeline/theory.py:1220}) marks which declared function plays the special
\code{phi\_concrete} role (it defaults to \code{'phi'}).

\begin{defn}
A SageMath \textbf{formal} (or \emph{undefined}) function is created by
\code{function('phi\_1')(v)}: an \emph{unapplied} symbolic function evaluated at
\code{v}, which Sage keeps as an opaque node that \code{taylor()} can
differentiate \emph{formally}, without ever needing a concrete body.
\end{defn}

The crucial design choice is that the action is authored with \code{phi} as a
\emph{formal} Sage function. When the field-theory core Taylor-expands the action,
Sage produces the formal-derivative symbols $\phi(\bar v),\,\phi'(\bar v),\,
\phi''(\bar v),\ldots$. The framework names these \code{phi0\_<i>},
\code{phi1\_<i>}, \code{phi2\_<i>}, \ldots (the \code{\_<i>} suffix is the 1-based
population index), and they become the \emph{vertex couplings} of the
diagrammatic expansion. The concrete numeric values of those derivatives are
substituted only later, by the \code{specializations} and
\code{mf\_bg\_conditions} lambdas.

\begin{gotcha}
If you inline the concrete body where the formal function should be, you
short-circuit the entire auto-Taylor machinery. The Hawkes template docstring is
explicit (\file{pipeline/theory\_templates.py:319}--\code{:345}): using the
concrete expression makes the auto-expander ``see \code{phi(0)=0} etc.\ as numeric
and never register symbols, leaving an uncancelled \code{(1,0)} tadpole in the
bigrade analysis.'' This is why \code{functions\_list()} returns the function as
the formal \code{lambda i, v: function(f'phi\_\{i+1\}')(v)}
(\file{pipeline/theory\_templates.py:344}).
\end{gotcha}

\section{Mean field: the two surfaces}

The saddle $\bar\varphi$ must be found. There are two authoring surfaces for the
mean-field equations: a modern one preferred for new theories, and a legacy one.

\subsection{The modern surface: \texttt{.equation()}}

\code{equation} (\file{pipeline/theory.py:1087}) declares one residual of a
\term{DAE} system: a \emph{differential-algebraic equation} system. Each call adds
one equation $\text{LHS} = \text{RHS}$, i.e.\ the residual
$\text{LHS}-\text{RHS}=0$. At the mean-field point the framework substitutes
$\Dt\to 0$ on the LHS and solves the resulting algebraic system by
\term{multi-start Newton}: seed the solver from many random starting points,
collect the converged roots, sort them, and pick the \code{fixed\_point\_index}-th.
The quartic OU saddle equation is declared as

\begin{lstlisting}
.equation(lhs='(Dt+mu)*x[i]', rhs='-eps*x[i]^3', population='pop')
\end{lstlisting}

The method does real validation (\file{pipeline/theory.py:1143}--\code{:1184}),
which is worth knowing because the error messages are precise:

\begin{enumerate}
  \item \code{lhs}/\code{rhs} must be non-empty strings.
  \item \code{Conv(...)} on \emph{either} side raises (regex
        \code{\textbackslash bConv\textbackslash s*\textbackslash(},
        case-insensitive): the DAE assumes a \emph{stationary} mean field, so any
        kernel convolution of a constant must be pre-collapsed by you (for a
        normalised kernel, that is just the constant).
  \item \code{Dt} in the \code{rhs} raises (regex
        \code{\textbackslash bDt\textbackslash b}): derivatives belong on the LHS
        only.
  \item The \code{kind} is classified \code{'differential'} if \code{Dt} appears
        as a word in the LHS, else \code{'algebraic'}.
\end{enumerate}

\code{stability\_analysis(True)} (\file{pipeline/theory.py:1187}) opts into a
linear-stability filter for bistable / differential theories: each converged root
is linearised (keeping $\Dt\to -\ii\omega$), the generalised eigenvalue problem
$(B - \ii\omega A)\,\delta x = 0$ is assembled, and \code{fixed\_point\_index} is
restricted to the linearly stable subset. It is off by default (vacuous for
all-algebraic theories).

\subsection{The legacy surface: \texttt{set\_mf\_equation}}

\code{set\_mf\_equation(saddle\_name, rhs\_text)} (\file{pipeline/theory.py:1069})
is the older per-saddle surface: you give the RHS string for each saddle directly,
as in \code{set\_mf\_equation('nstar', 'phi(vstar[i])')}. The framework
auto-emits the EOM-closure relation \code{phi0\_<i> = nstar[i]} so you need not
declare it.

The two surfaces are bridged automatically. At \code{build()} time,
\code{\_auto\_populate\_mf\_eqs\_from\_equations} (\file{pipeline/theory.py:1276})
walks each \code{.equation(...)} you declared, substitutes $\Dt\to 0$ in the LHS,
finds the single state-variable reference, infers the saddle name as
\code{<var>star}, and textually rewrites every state-variable name in the RHS to
its \code{*star} counterpart (using word-boundary regexes so \code{Em} is not hit
when rewriting \code{E}). It is a no-op if you used the legacy surface directly.
This bridge is why a theory that declares only \code{.equation(...)} still
populates the legacy machinery the field-theory sanity check depends on.

\section{Kernels}

A \term{kernel} $g(t)$ is a convolution filter (e.g.\ an exponential synapse
$g(t) = (1/\tau_g)\,\ee^{-t/\tau_g}\,\Theta(t)$). Diagram propagators live in
frequency space, so the pipeline needs the kernel's Fourier image
$\hat g(\omega)$. For the exponential synapse,
\[
  \hat g(\omega) \;=\; \frac{1}{1 + \ii\,\omega\,\tau_g}.
\]

\code{define\_kernel} (\file{pipeline/theory.py:370}) is the text surface. You must
supply at least one of two keyword strings (it raises if both are \code{None}):

\begin{description}
  \item[\code{freq\_image}] the frequency image directly, e.g.\
        \code{'1/(1 + I*omega*tau\_g)'}. This is the fast path --- no transform is
        attempted.
  \item[\code{time\_expr}] the time-domain $g(t)$, e.g.\
        \code{'(1/tau\_g)*exp(-t/tau\_g)*heaviside(t)'}. The compiler then
        \emph{symbolically} Fourier-transforms it.
\end{description}

The transform is done by \code{make\_kernel\_ft\_image\_lambda}
(\file{pipeline/theory\_compiler.py:1493}), which calls
\code{msrjd.core.field\_theory.fourier\_transform} --- a wrapper around Sage's
\code{integrate} computing $\int g(t)\,\ee^{-\ii\omega t}\,\dd t$ (the
$\ee^{-\ii\omega t}$ convention, consistent with the exponential-synapse image
$\hat g(\omega)=1/(1+\ii\omega\tau_g)$ below). The text may
reference \code{i}/\code{j} for vector/matrix kernels, e.g.\
\code{'1/(1 + I*omega*tau\_g[i, j])'}.

\begin{gotcha}
The symbolic Fourier transform is best-effort. Sage handles
\code{heaviside(t)*exp(-t/tau)} for typical neural kernels, but harder forms fail;
\code{make\_kernel\_ft\_image\_lambda} then raises a \code{ValueError} telling you
to supply \code{freq\_image} directly (\file{pipeline/theory\_compiler.py:1530}).
When in doubt, give the frequency image.
\end{gotcha}

\section{Noise beyond white Gaussian}

White Gaussian noise is the $-D\,\tilde\varphi^{2}$ term and needs no special
machinery. Two generalisations have dedicated surfaces.

\subsection{Higher cumulants: \texttt{declare\_cgf\_term}}

\code{declare\_cgf\_term} (\file{pipeline/theory.py:461}) declares one row of a
\term{CGF} --- cumulant generating functional --- term: a non-closed-form noise
cumulant. The key arguments are \code{order} (the cumulant order: 2 for
$\kappa^{(2)}$, 3 for $\kappa^{(3)}$, \ldots), a Sage-syntax \code{coefficient},
an optional time-domain \code{kernel} factor, and either \code{response\_field}
(legacy single field) or \code{response\_legs} (a per-leg list for cross-field
cumulants). Rows sharing a \code{name} and \code{order} are summed into one
cumulant.

\subsection{Coloured noise and the Markovian embedding}

\begin{defn}
\textbf{White} noise is $\delta$-correlated,
$\avg{\eta\,\eta'}\propto\delta(\tau)$; \textbf{coloured} noise has a finite
correlation time, e.g.\ $\avg{\eta\,\eta'}\propto c\,\ee^{-|\tau|/\tau_c}$ (a
single Lorentzian in frequency). Coloured noise makes the loop integrals
expensive --- they would hang in \code{scipy.nquad} at loop order $\ge 1$.
\end{defn}

The standard cure is a \term{Markovian embedding}: introduce an auxiliary
Ornstein--Uhlenbeck field driven by \emph{white} noise whose stationary
autocorrelation is exactly the desired exponential, and couple the original field
to it linearly. \code{markovianize(True)} (\file{pipeline/theory.py:429}, default
on) toggles this. When on, \code{build()} walks the declared CGF terms and
rewrites every row whose kernel matches the \code{c*exp(-abs(tau)/tauc)} template
into a white-noise auxiliary field plus a linear filter, \emph{before} the
compiler runs. The rewrite itself lives in
\file{pipeline/colored\_to\_markovian.py}; the builder only flips the switch.

\section{Templates}

For recurring model families there are pre-canned templates in
\file{pipeline/theory\_templates.py}, so you do not retype a standard action.

\begin{description}
  \item[\code{HawkesAction}] (\file{pipeline/theory\_templates.py:69}) --- the
        canonical 2-population Hawkes action generator. Its methods each return a
        lambda: \code{action()} (the full \msrjd{} action),
        \code{phi\_concrete()}, \code{specializations()},
        \code{mf\_bg\_conditions()}, \code{mf\_equations()},
        \code{mf\_substitutions()}, and
        \code{functions\_list()} (the formal $\phi_i$).
  \item[\code{ExpSynapticKernel}] (\file{pipeline/theory\_templates.py:353}) ---
        the exponential synapse with image
        $\hat g(\omega)=1/(1+\ii\omega\tau_g)$.
  \item[\code{DeltaKernel}] (\file{pipeline/theory\_templates.py:374}) --- an
        instantaneous $g(t)=\delta(t)$. Its \code{kernel\_ft\_image()} returns
        \code{None}; its \code{extra\_specializations()} returns
        \code{lambda ns: \{ns.g: ns.delta\_D\}} so the kernel symbol is replaced
        by the Dirac $\delta$ at propagator-construction time.
  \item[\code{GTaSNoise}] (\file{pipeline/theory\_templates.py:400}) --- a
        Bernoulli + Gaussian external-rate process supplying $\kappa^{(2)}$,
        $\kappa^{(3)}$, $\kappa^{(4)}$.
\end{description}

\code{use\_action\_template} (\file{pipeline/theory.py:665}) wires a template in by
pulling its seven companion lambdas off it (\code{action}, \code{phi\_concrete},
\code{specializations}, \code{mf\_bg\_conditions}, \code{mf\_equations},
\code{mf\_substitutions}, \code{functions\_list}). \code{add\_gtas\_noise}
(\file{pipeline/theory.py:591}) is a heavier convenience: it auto-declares the
GTaS fields and parameters and re-emits the action with the noise hook installed.

\begin{gotcha}
\code{add\_gtas\_noise} declares the response field \code{mt} with
\code{auto\_response=False} \emph{on purpose}. Leaving auto-response on would
declare \code{mt} twice, which produces a duplicate \code{mt1} generator in the
bigrade \code{PolynomialRing} and the hard error \code{ValueError: variable name
'mt1' appears more than once} (\file{pipeline/theory.py:606}).
\end{gotcha}

\section{Spatial theories and the operator IR}

For a spatial field $\varphi(x,t)$ the action contains spatial differential
operators. On a diagram leg of wavevector $k$ and frequency $\omega$ their Fourier
images --- the \term{form factors} the spatial integrator attaches to each leg ---
are
\[
  \Lap \to -k^{2}, \qquad \Dt \to -\ii\omega, \qquad \partial_{x_i} \to \ii k_i,
\]
and compositions multiply ($\nabla^{4}\to k^{4}$). There are two authoring styles.

\subsection{Plain (v1): the Laplacian as a bare symbol}

In a spatial theory \code{build()} registers a \code{Laplacian} operator symbol
that is used multiplicatively, exactly like \code{Dt}. A reaction--diffusion
action with a $g\,p^{2}$ vertex reads
\code{pt*(Dt(p) + mu*p - DD*Lap(p) + g*p\^{}2) - T*pt\^{}2}. This style is enough
whenever the derivative acts on a single leg.

\subsection{Operator IR (v2): \texttt{.operator\_ir()}}

\begin{defn}
The \textbf{operator IR} (intermediate representation) makes \code{Lap(phi)},
\code{Dt(phi)}, \code{Dx(phi, i)} \emph{argument-binding calls} rather than bare
multiplicative symbols. Turn it on with \code{.operator\_ir()}
(\file{pipeline/theory.py:1053}). It is required for
\emph{derivative-inside-nonlinearity} vertices, where a bare multiplicative symbol
could not say \emph{which} leg the derivative acts on.
\end{defn}

The motivating example is KPZ: $(\partial_x h)^{2}$ is \emph{not} the same vertex
as $\partial_x(h^{2})$, and only a call that binds its argument can capture the
difference. The KPZ action (\file{theories/kpz\_1d.theory.py}) is

\begin{lstlisting}
.set_action_text('ht*(Dt(h) + mu*h - D*Lap(h) - (lam/2)*Dx(h, 0)^2) - T*ht^2')
.operator_ir()
\end{lstlisting}

The lowering machinery lives in \file{pipeline/spatial\_operator\_ir.py}, driven
from the compiler's \code{\_lower\_operator\_ir\_action}
(\file{pipeline/theory\_compiler.py:720}). It applies operator linearity,
annihilates operators acting on the homogeneous saddle (\code{kill\_means}),
replaces each atomic \code{Op(fluctuation)} with a fresh ring generator (the
$u=\delta\varphi,\,v=\nabla^{2}\delta\varphi$ trick), and classifies the
generators into \term{bilinear} (folded into the propagator) versus \term{vertex}
(carrying per-leg form factors).

\begin{gotcha}
A derivative \emph{vertex} is restricted to base field-degree 1 (\emph{per-leg},
like KPZ $(\partial\varphi)^{2}$) or 2 (\emph{composite}, like Model~B
$\nabla^{2}(\varphi^{2})$ or Burgers $\partial_x(\varphi^{2})$); any other degree
raises \code{NotImplementedError} (\file{pipeline/theory\_compiler.py:784}). A
bilinear \code{Dx} on a transverse axis ($\ne 0$ for $d\ge 2$) also raises
(\code{:837}).
\end{gotcha}

\subsection{Boundary, initial conditions, and coupled diffusion}

\code{boundary(mode, ...)} (\file{pipeline/theory.py:148}) accepts
\code{'infinite'} or \code{'periodic'} (Dirichlet/Neumann/Robin are explicitly
v2); periodic requires a \code{length=}. \code{initial(mode)}
(\file{pipeline/theory.py:182}) supports only \code{'stationary'} in v1. For
coupled spatial theories with \emph{unequal} diffusion constants there is a
\term{Dyson--Duhamel} dressing of the retarded edges: \code{dyson(mode, order)}
(\file{pipeline/theory.py:206}, with sugar \code{dyson\_order}) sets the
fixed-order truncation, and \code{reference\_diffusion(D0)}
(\file{pipeline/theory.py:276}) sets the scalar reference $D_0$ in the split
$\mathcal{D} = D_0\,I + \hat{\mathcal{D}}$. Declaring any of these on a non-spatial
theory is a hard error.

\section{How text becomes lambdas}\label{sec:textcompile}

This is the machinery section-within-the-chapter: how a string like
\code{'sum(xt[i]*((Dt+mu)*x[i] + ...) ...)'} becomes a callable that returns an
\code{SR} expression. It all happens in \file{pipeline/theory\_compiler.py}, driven
from \code{\_compile\_text\_declarations} (\file{pipeline/theory.py:1367}).

\subsection{The preparse-then-eval core}

User text uses the mathematician's \code{\^{}} for powers and bare integer
literals. Sage normally rewrites such source with a \term{preparser} before
running it (most importantly \code{\^{}}~$\to$~\code{**}, and lifting integers to
Sage \code{Integer} objects). The compiler invokes that preparser by hand in
\code{\_safe\_eval} (\file{pipeline/theory\_compiler.py:682}):

\begin{lstlisting}
import sage.all as _sage_all
from sage.repl.preparse import preparse
text = _normalize_expr_text(text)
parsed = preparse(text)
eval_globals = {**_sage_all.__dict__, **locals_dict}
return eval(parsed, eval_globals)
\end{lstlisting}

Two subtleties are spelled out in the docstring (lines~682--701):

\begin{itemize}
  \item \textbf{Why preparse:} it performs \code{\^{}}~$\to$~\code{**} (so
        \code{a*v\^{}2} means $a\,v^{2}$, not the Python bitwise XOR) and the
        integer lift.
  \item \textbf{Why globals, not locals:} the namespace is passed as \code{eval}'s
        \emph{globals}, because Python's \code{eval} only feeds its \code{locals}
        argument to the \emph{outermost} scope. A generator expression inside
        \code{sum(... for j in pop)} opens an inner scope that would \emph{not}
        see \code{locals}; putting everything in globals fixes the comprehension
        scoping.
\end{itemize}

On any failure \code{\_safe\_eval} raises a rich \code{ValueError} quoting the
text, the preparsed form, the original error, and the sorted list of available
symbols --- the single most useful debugging affordance in the subsystem.

\begin{gotcha}
A consequence of globals-not-locals: user names \emph{shadow} Sage builtins. A
parameter literally named \code{I} would shadow the imaginary unit. This is
documented as intentional (\file{pipeline/theory\_compiler.py:697}) but is a
foot-gun.
\end{gotcha}

\subsection{The wrapper classes}

Before \code{\_safe\_eval} runs, the compiler builds a namespace dict mapping each
name you used to an object that makes the text behave. These wrapper objects live
only transiently inside that dict --- they never appear in the model dict --- but
they are the heart of the DSL:

\begin{description}
  \item[\code{\_FullPhysicalField}] (\file{pipeline/theory\_compiler.py:206}) ---
        exposes \code{n[i] = nstar[i] + dn[i]}, so writing \code{nt[i] * n[i]} is
        shorthand for \code{nt[i] * (nstar[i] + dn[i])}. Its
        \code{\_\_getitem\_\_} returns \code{saddle[i] + fluct[i]}.
  \item[\code{\_FieldScalar}] (\file{pipeline/theory\_compiler.py:268}) --- wraps a
        list of \code{SR} variables so a size-1 field supports \emph{both} bare
        scalar arithmetic (\code{xt * x}) and indexed access
        (\code{xt[0] * x[0]}); size $>1$ raises a clear \code{TypeError}.
  \item[\code{\_MatrixView}] (\file{pipeline/theory\_compiler.py:316}) --- lets a
        matrix parameter accept both \code{w[i, j]} and \code{w[i][j]}.
  \item[\code{\_IndexedFormalFunction}] (\file{pipeline/theory\_compiler.py:73})
        --- turns \code{phi[i](dv[i])} into \code{function('phi\_<i+1>')(dv[i])},
        the formal symbol the auto-Taylor pass needs.
\end{description}

\subsection{The unwrap dance}

There is one Sage wrinkle these wrappers force. Sage's strict argument coercion
inside \code{BuiltinFunction.\_\_call\_\_} (used by \code{exp}, \code{log}, and
formal functions) \emph{ignores} a Python object's \code{\_sage\_()} method. So
\code{exp(xt)} in scalar mode, where \code{xt} is a \code{\_FieldScalar}, raises
\code{no canonical coercion from \_FieldScalar to SR}. The fix is a pair of
helpers --- \code{\_unwrap\_field\_arg} (\file{pipeline/theory\_compiler.py:49})
and \code{\_unwrap\_builtin} (\file{pipeline/theory\_compiler.py:342}) --- that
unwrap a size-1 field wrapper to its bare \code{SR} expression \emph{before} it
reaches Sage's coercion. Every builtin exposed to user text is wrapped this way in
\code{\_builtin\_namespace} (\file{pipeline/theory\_compiler.py:359}); any new
builtin must be wrapped the same way or scalar-mode calls will break.

\subsection{The \texttt{ns} object the lambdas receive}

\begin{defn}
The lambdas all take one argument, \code{ns} --- a \code{NamespaceForFields}
object \emph{built by the consumer} (\code{msrjd.core.field\_theory}), not by this
subsystem. It carries the symbolic variables as attributes the lambdas read:
\code{ns.<fieldname>} (lists of \code{SR} vars), \code{ns.<paramname>},
\code{ns.<kernelname>}, \code{ns.pop}, \code{ns.Dt}, \code{ns.Laplacian}, and
several private bookkeeping attributes (\code{ns.\_deriv\_rename\_subs},
\code{ns.\_all\_field\_sr\_vars}, \code{ns.\_operator\_ir}, \code{ns.delta\_D}).
\end{defn}

So \code{make\_action\_lambda} (\file{pipeline/theory\_compiler.py:851}) does not
return an expression --- it returns a \emph{closure} \code{\_action(ns)} that,
when later called with a concrete namespace, builds the eval namespace from
\code{ns}, runs \code{\_safe\_eval} on your text, optionally lowers the operator
IR, kills any ``operator applied to a constant saddle'' terms (e.g.\
$\Dt\cdot\bar v = 0$, $\Lap\cdot\bar\varphi = 0$), and returns the symbolic action.

\begin{gotcha}
The ``operator-times-saddle $\to 0$'' kill rule is \emph{duplicated}: it appears
in both \code{make\_action\_lambda} and \code{make\_mf\_bg\_conditions\_lambda},
and they walk the expression tree independently. The comment at
\file{pipeline/theory\_compiler.py:1378} is explicit that the two must stay in
sync: ``a missed operator here silently leaves a non-zero saddle term and the MF
solver fails to converge.'' It is a fragile, manually-mirrored invariant.
\end{gotcha}

\section{The \texttt{build()} method and the model dict}

\code{build()} (\file{pipeline/theory.py:1688}) is the terminal call. Its sequence
is:

\begin{enumerate}
  \item \textbf{Scalar-mode autopop} --- if no populations were declared but
        physical fields exist (and \code{n\_populations <= 1}), append a single
        size-1 population \code{pop} and bind every field to it.
  \item \textbf{Markovianize} --- if there are CGF terms, rewrite the spec via
        \code{markovianize\_spec} \emph{before} compiling.
  \item \textbf{Compile text declarations} ---
        \code{\_compile\_text\_declarations()} turns every text string into a
        lambda.
  \item \textbf{Validate required hooks} --- \code{action} is always required;
        \code{phi\_concrete} is required only when a single-argument function
        exists; \code{kernel\_ft\_image} is optional. A missing hook raises a
        \code{ValueError} listing what is absent.
  \item \textbf{Reserved-name validation} (see below).
  \item \textbf{Assemble and return} the model dict.
\end{enumerate}

\begin{gotcha}
Scalar-mode autopop is gated on \code{n\_populations <= 1}
(\file{pipeline/theory.py:1697}). Firing it for a legacy 2-population Hawkes model
would clobber the population count to 1, drop the second population's response
fields, and zero every cross-population correlator. Such theories keep
\code{populations} empty and use the flat legacy index path instead.
\end{gotcha}

\subsection{Reserved names: a scoped hard error}

Some names are owned by the framework's symbol machinery; reusing one as a field,
parameter, function, or kernel name collides --- often \emph{silently} --- with a
Fourier-transform, operator, or spatial-coordinate symbol. \code{build()}
(\file{pipeline/theory.py:1850}) checks this and raises with a rename hint. The
reserved set is \emph{scoped}:

\begin{center}
\begin{tabular}{@{}lll@{}}
\toprule
\textbf{Names} & \textbf{Reserved in} & \textbf{Why} \\
\midrule
\code{t}, \code{omega} & all theories & Fourier-transform variables \\
\code{Dt}, \code{delta\_D}, \code{delta\_Dp} & all theories & operator symbols \\
\code{k} & spatial theories & wavevector (\code{SR.var('k')}) \\
\code{Laplacian} & spatial theories & diffusion operator \\
\code{x}, \code{y}, \code{z} & spatial theories & spatial coordinates / output axis \\
\bottomrule
\end{tabular}
\end{center}

\begin{gotcha}
Note the irony: \code{x} is the conventional default field name in a
\emph{temporal} theory (as in the quartic OU model), but it is reserved in a
\emph{spatial} theory, where it is the spatial coordinate. The comment at
\file{pipeline/theory.py:1842} flags that a parameter colliding with \code{k}
would be a \emph{silent wrong-mass bug} (because \file{heat\_kernel.py} uses
\code{SR.var('k')} for the wavevector), which is precisely why \code{k} is a hard
error in spatial theories. The build raises with a rename hint such as
\code{x}~$\to$~\code{phi}.
\end{gotcha}

\subsection{The keys of the model dict}

\code{build()} assembles the dict at \file{pipeline/theory.py:1943}. The
\emph{always-present} keys are:

\begin{lstlisting}[style=console]
name, populations, iteration_saddles, index_sets, response_fields,
physical_fields, parameters, kernels, equations, operators, functions,
mf_substitutions, phi_concrete, action, naming_convention,
stability_analysis, operator_ir
\end{lstlisting}

Conditionally attached lambda keys (present only when the corresponding text was
declared): \code{mf\_bg\_conditions\_action}, \code{mf\_bg\_conditions},
\code{mf\_equations}, \code{kernel\_ft\_image}, \code{kernel\_td\_image},
\code{specializations}, \code{correlated\_noises}. Spatial-only keys (present only
when a field is spatial): \code{spatial}, \code{boundary}, \code{initial}.

\begin{gotcha}
The mean-field conditions are stored \emph{twice} under two keys, and the reason
is subtle (\file{pipeline/theory.py:1974}--\code{:1986}).
\code{mf\_bg\_conditions\_action} is the closure-baked version, pre-substituted so
the \code{(1,0)} tadpole cancels symbolically --- it is what
\code{FieldTheory.expand} uses. \code{mf\_bg\_conditions} is the solver-friendly
version that keeps raw \code{nstar} so the numerical iteration can run. If only
one exists (the legacy template path), it is used for both.
\end{gotcha}

The \code{naming\_convention} key deserves a look, because it is how the rest of
the pipeline translates between the natural letters you typed and the internal
names:

\begin{lstlisting}
{
  'fluctuation_fields': {'n': 'dn', 'v': 'dv'},     # natural -> internal
  'mean_field_saddles': {'n': 'nstar', 'v': 'vstar'},  # natural -> internal
  'mf_parameters':      ['nstar', 'vstar'],          # internal names
}
\end{lstlisting}

\section{End to end: who consumes the dict}

The model dict is read by \code{compute\_cumulants}
(\file{pipeline/compute.py:108}), the function the quickstart chapter drove
through \code{dd.run}. Its docstring lists the \emph{required} hooks ---
\code{response\_fields}, \code{physical\_fields}, \code{parameters},
\code{kernels}, \code{operators}, \code{functions}, \code{mf\_substitutions},
\code{mf\_bg\_conditions}, \code{specializations}, \code{kernel\_ft\_image},
\code{phi\_concrete}, \code{mf\_equations}, \code{action} (with
\code{correlated\_noises} optional) --- and the builder produces exactly these.
The numeric parameter values come not from the dict but from the theory file's
\code{DEFAULT\_FUNDAMENTAL} constant, and the run recommendations from its
\code{METADATA}, as the previous chapter showed.

The whole flow, from the file you write to the numbers downstream, is:

\begin{center}
\begin{tikzpicture}[
    node distance=9mm and 16mm,
    box/.style={draw=noteframe, fill=notebg, rounded corners=2pt,
                inner sep=5pt, align=center, font=\small},
    op/.style={font=\footnotesize\itshape, align=center},
    >={Stealth[length=2.4mm]}]
  \node[box] (file) {\code{theories/<name>.theory.py}\\\code{build()} function};
  \node[box, right=of file] (builder) {\code{TheoryBuilder}\\(accumulates\\declarations)};
  \node[box, below=12mm of builder] (model) {model dict\\(numbers + Sage\\lambdas)};
  \node[box, left=of model] (compute) {\code{compute\_cumulants}\\(enumerate $\to$\\integrate)};
  \draw[->] (file) -- node[op, above]{text} (builder);
  \draw[->] (builder) -- node[op, right]{\code{.build()}} (model);
  \draw[->] (model) -- node[op, above]{required\\hooks} (compute);
\end{tikzpicture}
\end{center}

\section{Two complete examples}

The two theory files behind this and the previous chapter are short enough to read
whole. The temporal one (\file{theories/ou\_quartic.theory.py}) is the quartic OU
model of \eqref{eq:ouquartic-here}:

\begin{lstlisting}
from pipeline.theory import TemporalTheoryBuilder

def build():
    return (
        TemporalTheoryBuilder('OU Quartic (white noise)')
        .population('pop', size=1)
        .physical_field('x', population='pop', description='variable')
        .parameter('mu', default=1.0, domain='positive')
        .parameter('eps', default=0.02, domain='positive')
        .parameter('D', default=1.0, domain='positive')
        .set_action_text('sum(xt[i]*((Dt+mu)*x[i] + eps*x[i]^3) - D*xt[i]^2 for i in pop)')
        .equation(lhs='(Dt+mu)*x[i]', rhs='-eps*x[i]^3', population='pop')
        .build()
    )
\end{lstlisting}

Reading it as a chain: \code{.physical\_field('x')} declares the fluctuation
\code{dx}, the response \code{xt}, and the saddle \code{xstar}; the three
\code{.parameter(...)} calls add $\mu,\varepsilon,D$ (all positive); the action
text is \eqref{eq:ouquartic-here} verbatim; and \code{.equation(...)} is the
saddle condition $(\Dt+\mu)x = -\varepsilon x^{3}$. At \code{.build()},
scalar-mode autopop is skipped (a \code{pop} population was declared), the
equation is bridged into the legacy mean-field surface, and every text string is
compiled to a lambda.

The spatial one (\file{theories/kpz\_1d.theory.py}) is 1D KPZ, and shows the
operator IR in anger:

\begin{lstlisting}
from pipeline.theory import SpatialTheoryBuilder

def build():
    return (
        SpatialTheoryBuilder('1D KPZ (per-leg gradient vertex)', n_populations=0)
        .physical_field('h', spatial_dim=1, description='interface height')
        .parameter('mu',  default=1.0, domain='positive')
        .parameter('D',   default=1.0, domain='positive')
        .parameter('lam', default=0.3, domain='real')
        .parameter('T',   default=1.0, domain='positive')
        .equation(lhs='(Dt + mu - D*Laplacian)*h', rhs='0')
        .set_action_text('ht*(Dt(h) + mu*h - D*Lap(h) - (lam/2)*Dx(h, 0)^2) - T*ht^2')
        .operator_ir()
        .boundary('infinite')
        .initial('stationary')
        .build()
    )
\end{lstlisting}

Here \code{spatial\_dim=1} marks \code{h} as a 1D field (so \code{build()} will
register a \code{Laplacian} operator and emit a \code{spatial} block); the action
uses the binding-call form \code{Lap(h)}, \code{Dt(h)}, \code{Dx(h, 0)} because
the KPZ vertex $(\partial_x h)^{2}$ is a per-leg derivative vertex that the plain
multiplicative style cannot express; \code{.operator\_ir()} switches on the
lowering; and \code{.boundary('infinite')}, \code{.initial('stationary')} pin the
spatial setup. Note that the field is named \code{h}, not \code{x} --- because in
a spatial theory \code{x} is the reserved spatial coordinate.

\section{Where to go next}

You now know how the model dict is born. The next chapter, on theory files, covers
the \code{DEFAULT\_FUNDAMENTAL} and \code{METADATA} constants that sit alongside
\code{build()}, and a later chapter develops the coloured/Markovian noise embedding
sketched here. From the other side, the field-theory core chapter picks up the
model dict and starts the seven-phase pipeline you watched run in the quickstart:
it is the first consumer of the \code{action} lambda this subsystem so carefully
compiled.
